(1) 一维粒子质量 \(m\)，势能 \(V(x)=-V_0\)（\(|x|<a\)）、\(V(x)=0\)（\(|x|\ge a\)），其中 \(V_0,a>0\)。对束缚能量 \(-V_0<E<0\) 定义阱内外波数，并写出波函数基本形式。
(2) 利用宇称及 \(x=\pm a\) 的连续性，导出偶态和奇态各自的量子化条件。
(3) 证明任意浅的该吸引势阱至少具有一个偶束缚态，并求出现第一个奇束缚态的阈值。
(4) 对偶态写出归一化系数，讨论束缚能量接近零时波函数在阱外的空间尺度。
